\manualchapter{Patches and troubleshooting}{sec:operation}
\subsection{A three-note chord}
Set frequency offsets to 0, $4/12$ and $7/12$ octaves. Patch one pitch source
to the first V/OCT and one envelope to its Level; leave the corresponding
later inputs empty. Monitor the third output for a major triad. To sequence
the third note independently, patch its V/OCT; that replaces the inherited
pitch without changing the other two.

For a changing timbre, patch an audio oscillator into the first FM and raise
the depth from zero. FM is exponential, so strong modulation changes perceived
tuning. Use a small depth when a stable note center matters.

For silence, check Level and its CV. For a missing voice in the final mix,
check whether an earlier output is patched. For distortion, lower all three
levels before looking for a fault in the downstream mixer.

\subsection{Cable channels, outputs and saved patches}
The widest input selects 1--16 polyphonic chip instances. Voice columns are
oscillators within each instance, not separate cable channels. Inputs read
the matching channel; supply matching pitch and level channel counts rather
than expecting a mono envelope to repeat across a polyphonic pitch cable.

Each output adds the preceding output when that earlier socket is unpatched.
The accumulated signal clips at approximately $\pm5$\,V. To hear the complete
mix, use the last output and leave earlier outputs empty. Patching an earlier
output removes that signal from the following mix. Lower source levels if
the sum clips.

Rack saves the panel parameters. Reset initializes the chip; oscillator
phase is not preserved by a patch save. Pitch/FM update each sample, while
level and other register controls update every 16 samples.
